\section{Geometrica}
Modella il numero di prove necessarie per ottenere il primo successo in una sequenza di prove di bernoulli indipendenti, ciascuna con probabilità di successo $p$.
\[
\mathcal{G}(p)
\]

\begin{figure}[H]
\centering

\begin{minipage}{0.45\textwidth}
\centering
\begin{tikzpicture}
\begin{axis}[
    axis lines = middle,
    ylabel = {$\mathbb{P}(X=k)$},
    xlabel = {$k$},
    width = 0.90\textwidth,
    height = 0.90\textwidth,
    ymin = 0, ymax = 0.45,
    xmin = 0, xmax = 31,
    legend style={at={(axis cs: 1,0.43)}, anchor=north west, font=\tiny},
    legend cell align=left,
]

% --- Geometrica(p=0.3) ---
\addplot[
    only marks,
    black,
    mark=*,
    mark options={fill=black, scale=0.7}
]
coordinates {
(1,0.3)(2,0.21)(3,0.147)(4,0.1029)(5,0.07203)
(6,0.050421)(7,0.035295)(8,0.024706)(9,0.017294)
(10,0.012106)(11,0.008474)(12,0.005932)(13,0.004152)
(14,0.002907)(15,0.002035)(16,0.001425)
(17,0.000998)(18,0.000699)(19,0.000489)
(20,0.000342)(21,0.000240)(22,0.000168)
(23,0.000118)(24,0.000083)(25,0.000058)
(26,0.000041)(27,0.000029)(28,0.000020)
(29,0.000014)(30,0.000010)
};

% --- Geometrica(p=0.6) ---
\addplot[
    only marks,
    gray!60,
    mark=*,
    mark options={fill=gray!60, scale=0.7}
]
coordinates {
(1,0.6)(2,0.24)(3,0.096)(4,0.0384)(5,0.01536)
(6,0.006144)(7,0.002458)(8,0.000983)
(9,0.000393)(10,0.000157)(11,0.000063)
(12,0.000025)(13,0.000010)(14,0.000004)
(15,0.000002)(16,0.000001)
};

\end{axis}
\end{tikzpicture}
\end{minipage}
\hfill
\begin{minipage}{0.45\textwidth}
\centering

\begin{tikzpicture}
\begin{axis}[
    axis lines = middle,
    ylabel = {$F(k)$},
    xlabel = {$k$},
    width = 0.90\textwidth,
    height = 0.90\textwidth,
    ymin = 0, ymax = 1.1,
    xmin = 0, xmax = 31,
    legend style={at={(axis cs: 20,0.25)}, anchor=north west, font=\tiny},
    legend cell align=left,
]

% --- CDF Geometrica(p=0.3) ---
\addplot[
    jump mark left,
    black,
    mark=*,
    mark options={fill=black, scale=0.7}
]
coordinates{
(1,0.3)(2,0.51)(3,0.657)(4,0.7599)(5,0.83193)
(6,0.882351)(7,0.917646)(8,0.942352)
(9,0.959647)(10,0.971753)(11,0.980227)
(12,0.986159)(13,0.990311)(14,0.993218)
(15,0.995252)(16,0.996677)(17,0.997675)
(18,0.998374)(19,0.998863)(20,0.999205)
(21,0.999445)(22,0.999613)(23,0.999731)
(24,0.999814)(25,0.999872)(26,0.999913)
(27,0.999942)(28,0.999962)(29,0.999976)
(30,0.999983)
};
\addlegendentry{$p=0.3$}

% --- CDF Geometrica(p=0.6) ---
\addplot[
    jump mark left,
    gray!60,
    mark=*,
    mark options={fill=gray!60, scale=0.7}
]
coordinates{
(1,0.6)(2,0.84)(3,0.936)(4,0.9744)(5,0.98976)
(6,0.995904)(7,0.998362)(8,0.999345)
(9,0.999738)(10,0.999895)(11,0.999958)
(12,0.999983)(13,0.999993)(14,0.999997)
(15,0.999999)
};
\addlegendentry{$p=0.6$}

\end{axis}
\end{tikzpicture}

\end{minipage}
\end{figure}


\bigskip
\begin{table}[H]
\centering
\renewcommand{\arraystretch}{2}
\begin{tabular}{@{}ll@{}}
\toprule
\textbf{Supporto} 
& $\{1,2,3,\dots\}\subset\mathbb{Z}$ \\

\midrule

\textbf{Funzione di Densità} 
& $\mathbb{P}(X=k)=p(1-p)^{k-1}$,\quad $k=1,2,3,\dots$ \\

\midrule

\textbf{Funzione di Ripartizione}
& $F(k)=\mathbb{P}(X\le k)=1-(1-p)^k$ \\

\midrule

\textbf{Valore atteso} 
& $\mathbb{E}[X]=\dfrac{1}{p}$ \\

\midrule

\textbf{Varianza}
& $Var(X)=\dfrac{1-p}{p^2}$ \\

\midrule

\textbf{Funzione caratteristica}
& $\varphi_X(u)=\dfrac{pe^{iu}}{1-(1-p)e^{iu}}$\\

\bottomrule
\end{tabular}
\end{table}

\noindent La Geometrica è l’unica variabile aleatoria discreta che gode della
\emph{proprietà di assenza di memoria}.
Ciò significa che la probabilità di dover attendere ancora un certo numero di prove
non dipende dal numero di insuccessi già osservati.

\smallskip
\noindent 
Formalmente per ogni $m,n \in \mathbb{N}$ vale:
\[
\mathbb{P}(X > m+n \mid X > m)
=
\mathbb{P}(X > n).
\]


